\documentclass[11pt]{article}
\usepackage[margin=27mm]{geometry}
\usepackage{amsmath,amssymb,amsthm,hyperref}
\newtheorem{definition}{Definition}
\newtheorem{theorem}{Theorem}
\newcommand{\lean}[1]{}
\newcommand{\leanok}{}
\newcommand{\uses}[1]{}
\begin{document}
\begin{definition}[Exact excitation matrices]
\label{def:replica_excitation_projection}
For $\Omega\in\mathbb C^A$ and $B\subseteq\{1,\ldots,k\}$, set
\[R_B=\bigotimes_{j=1}^k T_j,\qquad
T_j=\begin{cases}I-|\Omega\rangle\langle\Omega|,&j\in B,\\
|\Omega\rangle\langle\Omega|,&j\notin B.\end{cases}\]
If $\|\Omega\|=1$, then each $R_B$ is an orthogonal projection.
The empty tensor product is the identity on $\mathbb C$.
\end{definition}
\input{blueprint/src/fragment/replica_excitation_symmetry.tex}
\begin{thebibliography}{1}
\bibitem{OpenAI2026AreaLaw} OpenAI, \emph{A two-dimensional area law from a
global spectral gap}, September 24, 2026.
\end{thebibliography}
\end{document}
